% GENERATED FILE. Edit canonical lessons, then run npm run generate:books.
% canonical-source-hash: fnv1a64:305ca4651ee8fcfa
% canonical-lessons: ZH-C343-yueguang, ZH-C343-duonian, ZH-C343-xinzhong, ZH-C343-zhufu, ZH-C343-fenxiang

\chapter{Far from Home}
\label{ch:zh-far-from-home}
\hlchaptermodality{343}

\begin{chapteropening}
\textbf{By the end of this chapter:} \emph{I can talk about moonlight, years apart, what is in my heart, good wishes and sharing.}

Complete the last lesson of chapter 343: I can talk about moonlight, years apart, what is in my heart, good wishes and sharing.
\end{chapteropening}

\section[yuèguāng]{yuèguāng — moonlight}

\label{lesson:ZH-C343-yueguang}



\begin{quote}
Before the new one: say the Mandarin for sweet osmanthus flowers, then the Mandarin for one's native place.
\end{quote}

\subsection*{You'll want to know: yuèguāng}
\emph{yuèguāng} — \textquotedblleft{}moonlight\textquotedblright{}.

\begin{quote}
\emph{jīntiān wǎnshang de yuèguāng hěn hǎo} — The moonlight is lovely tonight.
\end{quote}

Its characters come later; for now it is a word for the ear.

\begin{cousinweb}
\emph{yuèguāng} is built from \zh{月} 'the moon' and \emph{guāng} 'light', both inherited from Middle Chinese, \emph{ngjwot} and \emph{kwang} (Guangyun, in Baxter's transcription), which glosses the second 'bright'.
\end{cousinweb}

\begin{grammarlens}[title={What you've built}]
One more word for telling what happened.
\end{grammarlens}

\subsection*{Your turn}
\begin{itemize}
\raggedright
  \item \emph{Say it:} \emph{yuèguāng}
  \item \emph{Say it:} \emph{yuèguāng}, once more
  \item \emph{Say it:} \emph{jīntiān wǎnshang de yuèguāng hěn hǎo} — The moonlight is lovely tonight.
  \item \emph{Recall:} say the Mandarin for sweet osmanthus flowers, then the Mandarin for one's native place, then say \emph{yuèguāng} again
\end{itemize}

\begin{tcolorbox}[breakable,colback=teal!4,colframe=teal!35!black,title={Before you move on}]
What does yuèguāng mean? (\textquotedblleft{}Moonlight\textquotedblright{}.) Say it once more.
\end{tcolorbox}
\section[\texorpdfstring{\zh{多年} (duōnián)}{多年 (duōnián)}]{\zh{多年} (duōnián) — for many years}

\label{lesson:ZH-C343-duonian}



\begin{quote}
Before the new one: say the Mandarin for one's native place, then the Mandarin for bright moonlight.
\end{quote}

\subsection*{You'll want to know: \zh{多年}}
\textbf{\zh{多年}} — \emph{duōnián} — \textquotedblleft{}for many years\textquotedblright{}.

\begin{quote}
\textbf{\zh{我们多年没有一起过中秋了}} \emph{wǒmen duōnián méiyǒu yìqǐ guò Zhōngqiū le} — We haven't spent the moon festival together for years.
\end{quote}

\begin{cousinweb}
\textbf{\zh{多年}} (\emph{duōnián}) is built from \zh{多} 'many' and \zh{年} 'a year', both inherited from Middle Chinese, \emph{ta} and \emph{nen} (Guangyun, in Baxter's transcription), which glosses the first 'many'.
\end{cousinweb}

\begin{grammarlens}[title={What you've built}]
One more word for telling what happened.
\end{grammarlens}

\subsection*{Your turn}
\begin{itemize}
\raggedright
  \item \emph{Say it:} \emph{duōnián}
  \item \emph{Say it:} \emph{duōnián}, once more
  \item \emph{Say it:} \emph{wǒmen duōnián méiyǒu yìqǐ guò Zhōngqiū le} — We haven't spent the moon festival together for years.
  \item \emph{Recall:} say the Mandarin for one's native place, then the Mandarin for bright moonlight, then say \emph{duōnián} again
\end{itemize}

\begin{tcolorbox}[breakable,colback=teal!4,colframe=teal!35!black,title={Before you move on}]
What does \zh{多年} mean? (\textquotedblleft{}For many years\textquotedblright{}.) Say it once more.
\end{tcolorbox}
\section[\texorpdfstring{\zh{心中} (xīnzhōng)}{心中 (xīnzhōng)}]{\zh{心中} (xīnzhōng) — in one's heart}

\label{lesson:ZH-C343-xinzhong}



\begin{quote}
Before the new one: say the Mandarin for bright moonlight, then the Mandarin for for many years now.
\end{quote}

\subsection*{You'll want to know: \zh{心中}}
\textbf{\zh{心中}} — \emph{xīnzhōng} — \textquotedblleft{}in one's heart\textquotedblright{}.

\begin{quote}
\textbf{\zh{在我心中这是最好的中秋}} \emph{zài wǒ xīnzhōng zhè shì zuì hǎo de Zhōngqiū} — To my mind, this is the best moon festival ever.
\end{quote}

\begin{cousinweb}
\textbf{\zh{心中}} (\emph{xīnzhōng}) is built from \zh{心} 'the heart' and \zh{中} 'the middle', both inherited from Middle Chinese, \emph{sim} and \emph{trjuwng} (Guangyun, in Baxter's transcription).
\end{cousinweb}

\begin{grammarlens}[title={What you've built}]
One more word for telling what happened.
\end{grammarlens}

\subsection*{Your turn}
\begin{itemize}
\raggedright
  \item \emph{Say it:} \emph{xīnzhōng}
  \item \emph{Say it:} \emph{xīnzhōng}, once more
  \item \emph{Say it:} \emph{zài wǒ xīnzhōng zhè shì zuì hǎo de Zhōngqiū} — To my mind, this is the best moon festival ever.
  \item \emph{Recall:} say the Mandarin for bright moonlight, then the Mandarin for for many years now, then say \emph{xīnzhōng} again
\end{itemize}

\begin{tcolorbox}[breakable,colback=teal!4,colframe=teal!35!black,title={Before you move on}]
What does \zh{心中} mean? (\textquotedblleft{}In one's heart\textquotedblright{}.) Say it once more.
\end{tcolorbox}
\section[zhùfú]{zhùfú — blessing; good wishes}

\label{lesson:ZH-C343-zhufu}



\begin{quote}
Before the new one: say the Mandarin for for many years now, then the Mandarin for deep in one's heart.
\end{quote}

\subsection*{You'll want to know: zhùfú}
\emph{zhùfú} — \textquotedblleft{}blessing; good wishes\textquotedblright{}.

\begin{quote}
\emph{Zhōngqiū nà tiān wǒ shōudào le péngyou de zhùfú} — On the day of the moon festival I got good wishes from my friends.
\end{quote}

Its characters come later; for now it is a word for the ear.

\begin{cousinweb}
\emph{zhùfú} is built from \zh{祝} 'to bless' and \emph{fú} 'good fortune', both inherited from Middle Chinese, \emph{tsyuwH} and \emph{pjuwk} (Guangyun, in Baxter's transcription), which glosses the second 'virtue' and 'divine help'.
\end{cousinweb}

\begin{grammarlens}[title={What you've built}]
One more word for telling what happened.
\end{grammarlens}

\subsection*{Your turn}
\begin{itemize}
\raggedright
  \item \emph{Say it:} \emph{zhùfú}
  \item \emph{Say it:} \emph{zhùfú}, once more
  \item \emph{Say it:} \emph{Zhōngqiū nà tiān wǒ shōudào le péngyou de zhùfú} — On the day of the moon festival I got good wishes from my friends.
  \item \emph{Recall:} say the Mandarin for for many years now, then the Mandarin for deep in one's heart, then say \emph{zhùfú} again
\end{itemize}

\begin{tcolorbox}[breakable,colback=teal!4,colframe=teal!35!black,title={Before you move on}]
What does zhùfú mean? (\textquotedblleft{}Blessing; good wishes\textquotedblright{}.) Say it once more.
\end{tcolorbox}
\section[fēnxiǎng]{fēnxiǎng — to share}

\label{lesson:ZH-C343-fenxiang}



\begin{quote}
Before the new one: say the Mandarin for deep in one's heart, then the Mandarin for good wishes sent to someone.
\end{quote}

\subsection*{You'll want to know: fēnxiǎng}
\emph{fēnxiǎng} — \textquotedblleft{}to share\textquotedblright{}.

\begin{quote}
\emph{wǒ hé tóngshì fēnxiǎng le wǒ de yuèbǐng} — I shared my mooncakes with my colleagues.
\end{quote}

Its characters come later; for now it is a word for the ear.

\begin{cousinweb}
\emph{fēnxiǎng} is built from \emph{fēn} 'to divide' and \emph{xiǎng} 'to offer', both inherited from Middle Chinese, \emph{pjun} and \emph{xjangX} (Guangyun, in Baxter's transcription), which glosses the first 'to hand out' and lists the second under an older spelling, 'to offer up'.
\end{cousinweb}

\begin{grammarlens}[title={What you've built}]
One more word for telling what happened.
\end{grammarlens}

\subsection*{Your turn}
\begin{itemize}
\raggedright
  \item \emph{Say it:} \emph{fēnxiǎng}
  \item \emph{Say it:} \emph{fēnxiǎng}, once more
  \item \emph{Say it:} \emph{wǒ hé tóngshì fēnxiǎng le wǒ de yuèbǐng} — I shared my mooncakes with my colleagues.
  \item \emph{Recall:} say the Mandarin for deep in one's heart, then the Mandarin for good wishes sent to someone, then say \emph{fēnxiǎng} again
\end{itemize}

\begin{tcolorbox}[breakable,colback=teal!4,colframe=teal!35!black,title={Before you move on}]
What does fēnxiǎng mean? (\textquotedblleft{}To share\textquotedblright{}.) Say it once more.
\end{tcolorbox}
